\chapter{Discussão}
\label{cap:discussao}

\section{Respostas às questões de pesquisa}
\label{sec:respostas}

\textbf{RQ1 --- integração segura.} A arquitetura adotada --- IA em processo
separado, contrato de dados validado, intenções restritas e sanitização em
duas camadas independentes --- mostrou-se suficiente para que a IA
participasse da ficção (conversa, comentários sobre o estado do mundo, tom
variável) sem receber qualquer capacidade de alterar o jogo. Em todos os
cenários de falha exercitados pelos testes --- serviço fora do ar, resposta
inválida, modelo lento, JSON malformado, intenção desconhecida --- o pior
resultado observado foi uma resposta de \emph{fallback}: nunca um travamento,
nunca um comando executado.

\textbf{RQ2 --- viabilidade em hardware modesto.} O jogo manteve 60 quadros
por segundo (limitados por sincronismo vertical) em 1280$\times$720 na GPU
integrada, com o cenário final completo, e a conversa local respondeu entre
12,9~s e 22,0~s com um modelo de 8 bilhões de parâmetros. Essa ordem de
grandeza é adequada à interação conversacional por turnos --- o formato
adotado no jogo --- e inadequada a reações em tempo real, o que reforça a
decisão de manter a IA fora do laço de controle.

\textbf{RQ3 --- rastreabilidade.} A combinação de especificação prévia,
critérios de aceitação verificáveis e suíte automatizada permitiu conduzir um
desenvolvimento de autor único, com apoio de agentes de IA, mantendo cada
afirmação do projeto atrelada a um teste ou a uma medição. As seis etapas
fecharam com a suíte verde, e as pendências reais do protótipo (validação
manual, empacotamento, modelo-alvo) estão declaradas --- não maquiadas.

\section{Interpretação dos resultados}
\label{sec:interpretacao}

Três leituras merecem destaque. A primeira é que o valor educacional do
artefato não foi medido: o que este trabalho demonstra é a viabilidade
técnica de um instrutor ficcional sempre disponível, com custo de
infraestrutura zero (nenhum serviço pago, nenhuma chave de API) e degradação
segura. A segunda é que a latência de segundos molda o \emph{design} do
diálogo: perguntas e respostas funcionam bem; reações instantâneas, não. A
terceira é que o detalhamento visual tem custo gerenciável quando segue
convenções (escala, nomenclatura, colisão embutida): o cenário cresceu sem
comprometer o desempenho.

\section{Comparação com a literatura}
\label{sec:comparacao}

A comparação com os trabalhos relacionados é qualitativa, conforme declarado
no \autoref{cap:relacionados} e sintetizado no Quadro~\ref{quad:comparativo}.
O contraste central é de propósito: \emph{Generative Agents}
\cite{park2023generative} e \emph{Voyager} \cite{wang2023voyager} investigam o
que a autonomia produz; o Laboratório 404 investiga o que se preserva quando
a autonomia é removida. Esse desenho conversa diretamente com as limitações
elencadas no levantamento de Gallotta \emph{et al.}
\cite{gallotta2024largelanguage} --- confiabilidade, custo e latência ---,
convertendo-as em restrições explícitas de projeto. Não há, contudo,
experimento comparativo controlado entre os sistemas; uma comparação nesse
nível exigiria métricas comuns (por exemplo, qualidade percebida do diálogo,
custo por sessão, taxa de comportamento inesperado), o que este trabalho não
produziu [ESTUDO COMPARATIVO NÃO REALIZADO].

\section{Contribuições}
\label{sec:contribuicoes}

As contribuições deste trabalho são:

\begin{enumerate}
  \item \textbf{Um padrão de arquitetura de IA diegética com capacidade
  limitada}, com contrato validado nos dois lados, sanitização em duas
  camadas e \emph{fallback} obrigatório --- aplicável a outros jogos
  educacionais que queiram incorporar LLMs com previsibilidade;
  \item \textbf{um processo incremental rastreável}, com especificação,
  critérios de aceitação, matriz de rastreabilidade e evidências automatizadas
  e visuais, adequado a desenvolvimento de autor único com apoio de agentes de
  IA;
  \item \textbf{um artefato completo e aberto}: jogo, assets, adaptador,
  suíte de testes, capturas e vídeo \cite{lab404repo,lab404video};
  \item \textbf{medições de desempenho e latência} em hardware sem GPU
  dedicada, com condições e limitações declaradas; e
  \item \textbf{um registro de lições de engenharia e de ambiente} que afetam
  a reprodutibilidade de projetos semelhantes (\autoref{sec:licoes}).
\end{enumerate}

\section{Limitações}
\label{sec:limitacoes}

As limitações são deliberadamente explícitas:

\begin{itemize}
  \item \textbf{Sem estudo com usuários.} Não há avaliação de aprendizagem,
  usabilidade ou experiência; o valor educacional permanece como hipótese
  [ESTUDO DE USUÁRIO NÃO REALIZADO].
  \item \textbf{Amostra pequena de latência.} Três contextos, uma execução
  cada; sem repetição estatística, sem variação de carga e sem comparação
  entre modelos.
  \item \textbf{Modelo-alvo não executado.} As medições usaram um modelo maior
  e mais lento (\texttt{llama3.1:8b}); os resultados são um limite superior
  conservador, mas o alvo (\texttt{llama3.2:3b}) não foi medido.
  \item \textbf{Validações manuais pendentes.} Mouse com captura de cursor,
  joystick físico e percurso humano completo do início ao fim.
  \item \textbf{Empacotamento aberto.} Não há executável independente gerado
  no ambiente utilizado (ausência de modelos de exportação do motor).
  \item \textbf{Referências normativas a validar.} As edições vigentes das
  normas ABNT e da norma de CLP citadas precisam de conferência final
  [VALIDAR COM O AUTOR], assim como o modelo institucional de formatação.
\end{itemize}

\section{Lições de engenharia e reprodutibilidade}
\label{sec:licoes}

O desenvolvimento revelou quatro classes de obstáculos que outros projetos
semelhantes provavelmente encontrarão, e que foram documentados no repositório
\cite{lab404repo}:

\begin{enumerate}
  \item \textbf{Confinamento de pacote do motor.} A instalação do motor por
  pacote confinado (formato \emph{snap}, no Linux) bloqueava o acesso a
  dispositivos de entrada até a conexão explícita da permissão de joystick;
  o sintoma era o adaptador existir no sistema e não existir para o jogo. O
  diagnóstico foi possível porque o jogo registra em log os dispositivos
  detectados na inicialização;
  \item \textbf{Estado de áudio por aplicativo.} O servidor de áudio mantinha
  um estado de mudo \emph{por aplicativo}; o jogo aparecia no mixer com
  volume integral e permanecia inaudível. A correção (desmutar o fluxo
  específico) persistiu para as execuções seguintes --- uma armadilha comum
  de depuração, porque todo teste de sistema tocava normalmente;
  \item \textbf{Limites do formato de transmissão 3D.} O formato glTF não
  anima emissão de material nem transporta texto dependente de fonte; LEDs
  intermitentes viraram lógica de script e textos de tela viraram nós de
  texto no motor --- decisões que também simplificaram ajustes posteriores;
  \item \textbf{Formato de importação e suposições de código.} O cálculo de
  repetição de áudio assumia dados de amostra não comprimidos; com a
  compressão aplicada na importação, o laço ficou incorreto até a correção
  --- e o episódio virou uma verificação de regressão.
\end{enumerate}

\section{Implicações}
\label{sec:implicacoes}

Para o \textbf{ensino de automação}, o trabalho sugere que um laboratório
ficcional de baixo custo pode ocupar o espaço entre a teoria e a bancada
física, com a IA no papel de instrutor consultivo. Para a \textbf{pesquisa em
IA em jogos}, oferece um ponto de contraste: os mesmos modelos que rendem mais
quando recebem autonomia também entregam valor narrativo quando recebem
limites explícitos e verificáveis. Para a \textbf{prática de desenvolvimento},
indica que um projeto individual assistido por agentes de IA pode manter
rastreabilidade quando a especificação e os testes precedem a implementação
--- e quando cada promessa é verificável por um comando.

\section{Trabalhos futuros}
\label{sec:futuros}

\begin{enumerate}
  \item Executar e comparar o modelo-alvo \texttt{llama3.2:3b} com o modelo
  usado nas medições, registrando latências nas mesmas condições;
  \item coletar métricas de tempo de quadro por percentil, para caracterizar
  engasgos transitórios além da média limitada por sincronismo;
  \item concluir o empacotamento do executável e a demonstração sem
  intervenção do desenvolvedor (FR-033);
  \item realizar o percurso humano completo e, em seguida, um estudo
  controlado com estudantes, medindo aprendizagem percebida e usabilidade;
  \item ampliar o conteúdo didático --- por exemplo, cenários de falha
  adicionais com sinais de campo --- mantendo a mesma disciplina de
  especificação e evidências; e
  \item automatizar a execução da suíte em integração contínua, com registro
  público dos resultados a cada alteração.
\end{enumerate}

\section{Considerações do capítulo}
\label{sec:consideracoes_discussao}

O balanço é de um protótipo tecnicamente sólido e honestamente delimitado: o
que foi verificado está demonstrado com números e evidências; o que não foi
verificado está listado, com nome e sobrenome. Essa separação é, ela própria,
parte do resultado.
